% Options for packages loaded elsewhere
% Options for packages loaded elsewhere
\PassOptionsToPackage{unicode}{hyperref}
\PassOptionsToPackage{hyphens}{url}
\PassOptionsToPackage{dvipsnames,svgnames,x11names}{xcolor}
%
\documentclass[
]{article}
\usepackage{xcolor}
\usepackage[margin=1in]{geometry}
\usepackage{amsmath,amssymb}
\setcounter{secnumdepth}{5}
\usepackage{iftex}
\ifPDFTeX
  \usepackage[T1]{fontenc}
  \usepackage[utf8]{inputenc}
  \usepackage{textcomp} % provide euro and other symbols
\else % if luatex or xetex
  \usepackage{unicode-math} % this also loads fontspec
  \defaultfontfeatures{Scale=MatchLowercase}
  \defaultfontfeatures[\rmfamily]{Ligatures=TeX,Scale=1}
\fi
\usepackage{lmodern}
\ifPDFTeX\else
  % xetex/luatex font selection
\fi
% Use upquote if available, for straight quotes in verbatim environments
\IfFileExists{upquote.sty}{\usepackage{upquote}}{}
\IfFileExists{microtype.sty}{% use microtype if available
  \usepackage[]{microtype}
  \UseMicrotypeSet[protrusion]{basicmath} % disable protrusion for tt fonts
}{}
\makeatletter
\@ifundefined{KOMAClassName}{% if non-KOMA class
  \IfFileExists{parskip.sty}{%
    \usepackage{parskip}
  }{% else
    \setlength{\parindent}{0pt}
    \setlength{\parskip}{6pt plus 2pt minus 1pt}}
}{% if KOMA class
  \KOMAoptions{parskip=half}}
\makeatother
% Make \paragraph and \subparagraph free-standing
\makeatletter
\ifx\paragraph\undefined\else
  \let\oldparagraph\paragraph
  \renewcommand{\paragraph}{
    \@ifstar
      \xxxParagraphStar
      \xxxParagraphNoStar
  }
  \newcommand{\xxxParagraphStar}[1]{\oldparagraph*{#1}\mbox{}}
  \newcommand{\xxxParagraphNoStar}[1]{\oldparagraph{#1}\mbox{}}
\fi
\ifx\subparagraph\undefined\else
  \let\oldsubparagraph\subparagraph
  \renewcommand{\subparagraph}{
    \@ifstar
      \xxxSubParagraphStar
      \xxxSubParagraphNoStar
  }
  \newcommand{\xxxSubParagraphStar}[1]{\oldsubparagraph*{#1}\mbox{}}
  \newcommand{\xxxSubParagraphNoStar}[1]{\oldsubparagraph{#1}\mbox{}}
\fi
\makeatother

\usepackage{color}
\usepackage{fancyvrb}
\newcommand{\VerbBar}{|}
\newcommand{\VERB}{\Verb[commandchars=\\\{\}]}
\DefineVerbatimEnvironment{Highlighting}{Verbatim}{commandchars=\\\{\}}
% Add ',fontsize=\small' for more characters per line
\usepackage{framed}
\definecolor{shadecolor}{RGB}{241,243,245}
\newenvironment{Shaded}{\begin{snugshade}}{\end{snugshade}}
\newcommand{\AlertTok}[1]{\textcolor[rgb]{0.68,0.00,0.00}{#1}}
\newcommand{\AnnotationTok}[1]{\textcolor[rgb]{0.37,0.37,0.37}{#1}}
\newcommand{\AttributeTok}[1]{\textcolor[rgb]{0.40,0.45,0.13}{#1}}
\newcommand{\BaseNTok}[1]{\textcolor[rgb]{0.68,0.00,0.00}{#1}}
\newcommand{\BuiltInTok}[1]{\textcolor[rgb]{0.00,0.23,0.31}{#1}}
\newcommand{\CharTok}[1]{\textcolor[rgb]{0.13,0.47,0.30}{#1}}
\newcommand{\CommentTok}[1]{\textcolor[rgb]{0.37,0.37,0.37}{#1}}
\newcommand{\CommentVarTok}[1]{\textcolor[rgb]{0.37,0.37,0.37}{\textit{#1}}}
\newcommand{\ConstantTok}[1]{\textcolor[rgb]{0.56,0.35,0.01}{#1}}
\newcommand{\ControlFlowTok}[1]{\textcolor[rgb]{0.00,0.23,0.31}{\textbf{#1}}}
\newcommand{\DataTypeTok}[1]{\textcolor[rgb]{0.68,0.00,0.00}{#1}}
\newcommand{\DecValTok}[1]{\textcolor[rgb]{0.68,0.00,0.00}{#1}}
\newcommand{\DocumentationTok}[1]{\textcolor[rgb]{0.37,0.37,0.37}{\textit{#1}}}
\newcommand{\ErrorTok}[1]{\textcolor[rgb]{0.68,0.00,0.00}{#1}}
\newcommand{\ExtensionTok}[1]{\textcolor[rgb]{0.00,0.23,0.31}{#1}}
\newcommand{\FloatTok}[1]{\textcolor[rgb]{0.68,0.00,0.00}{#1}}
\newcommand{\FunctionTok}[1]{\textcolor[rgb]{0.28,0.35,0.67}{#1}}
\newcommand{\ImportTok}[1]{\textcolor[rgb]{0.00,0.46,0.62}{#1}}
\newcommand{\InformationTok}[1]{\textcolor[rgb]{0.37,0.37,0.37}{#1}}
\newcommand{\KeywordTok}[1]{\textcolor[rgb]{0.00,0.23,0.31}{\textbf{#1}}}
\newcommand{\NormalTok}[1]{\textcolor[rgb]{0.00,0.23,0.31}{#1}}
\newcommand{\OperatorTok}[1]{\textcolor[rgb]{0.37,0.37,0.37}{#1}}
\newcommand{\OtherTok}[1]{\textcolor[rgb]{0.00,0.23,0.31}{#1}}
\newcommand{\PreprocessorTok}[1]{\textcolor[rgb]{0.68,0.00,0.00}{#1}}
\newcommand{\RegionMarkerTok}[1]{\textcolor[rgb]{0.00,0.23,0.31}{#1}}
\newcommand{\SpecialCharTok}[1]{\textcolor[rgb]{0.37,0.37,0.37}{#1}}
\newcommand{\SpecialStringTok}[1]{\textcolor[rgb]{0.13,0.47,0.30}{#1}}
\newcommand{\StringTok}[1]{\textcolor[rgb]{0.13,0.47,0.30}{#1}}
\newcommand{\VariableTok}[1]{\textcolor[rgb]{0.07,0.07,0.07}{#1}}
\newcommand{\VerbatimStringTok}[1]{\textcolor[rgb]{0.13,0.47,0.30}{#1}}
\newcommand{\WarningTok}[1]{\textcolor[rgb]{0.37,0.37,0.37}{\textit{#1}}}

\usepackage{longtable,booktabs,array}
\newcounter{none} % for unnumbered tables
\usepackage{calc} % for calculating minipage widths
% Correct order of tables after \paragraph or \subparagraph
\usepackage{etoolbox}
\makeatletter
\patchcmd\longtable{\par}{\if@noskipsec\mbox{}\fi\par}{}{}
\makeatother
% Allow footnotes in longtable head/foot
\IfFileExists{footnotehyper.sty}{\usepackage{footnotehyper}}{\usepackage{footnote}}
\makesavenoteenv{longtable}
\usepackage{graphicx}
\makeatletter
\newsavebox\pandoc@box
\newcommand*\pandocbounded[1]{% scales image to fit in text height/width
  \sbox\pandoc@box{#1}%
  \Gscale@div\@tempa{\textheight}{\dimexpr\ht\pandoc@box+\dp\pandoc@box\relax}%
  \Gscale@div\@tempb{\linewidth}{\wd\pandoc@box}%
  \ifdim\@tempb\p@<\@tempa\p@\let\@tempa\@tempb\fi% select the smaller of both
  \ifdim\@tempa\p@<\p@\scalebox{\@tempa}{\usebox\pandoc@box}%
  \else\usebox{\pandoc@box}%
  \fi%
}
% Set default figure placement to htbp
\def\fps@figure{htbp}
\makeatother





\setlength{\emergencystretch}{3em} % prevent overfull lines

\providecommand{\tightlist}{%
  \setlength{\itemsep}{0pt}\setlength{\parskip}{0pt}}



 


\makeatletter
\@ifpackageloaded{caption}{}{\usepackage{caption}}
\AtBeginDocument{%
\ifdefined\contentsname
  \renewcommand*\contentsname{Table of contents}
\else
  \newcommand\contentsname{Table of contents}
\fi
\ifdefined\listfigurename
  \renewcommand*\listfigurename{List of Figures}
\else
  \newcommand\listfigurename{List of Figures}
\fi
\ifdefined\listtablename
  \renewcommand*\listtablename{List of Tables}
\else
  \newcommand\listtablename{List of Tables}
\fi
\ifdefined\figurename
  \renewcommand*\figurename{Figure}
\else
  \newcommand\figurename{Figure}
\fi
\ifdefined\tablename
  \renewcommand*\tablename{Table}
\else
  \newcommand\tablename{Table}
\fi
}
\@ifpackageloaded{float}{}{\usepackage{float}}
\floatstyle{ruled}
\@ifundefined{c@chapter}{\newfloat{codelisting}{h}{lop}}{\newfloat{codelisting}{h}{lop}[chapter]}
\floatname{codelisting}{Listing}
\newcommand*\listoflistings{\listof{codelisting}{List of Listings}}
\makeatother
\makeatletter
\makeatother
\makeatletter
\@ifpackageloaded{caption}{}{\usepackage{caption}}
\@ifpackageloaded{subcaption}{}{\usepackage{subcaption}}
\makeatother
\usepackage{bookmark}
\IfFileExists{xurl.sty}{\usepackage{xurl}}{} % add URL line breaks if available
\urlstyle{same}
% fallback for those not using the hyperref driver hyperxmp:
\makeatletter
\@ifundefined{xmpquote}{\newcommand{\xmpquote}[1]{#1}}{}
\makeatother
\hypersetup{
  colorlinks=true,
  linkcolor={blue},
  filecolor={Maroon},
  citecolor={Blue},
  urlcolor={Blue},
  pdfcreator={LaTeX via pandoc}}


\author{}
\date{}
\begin{document}

\renewcommand*\contentsname{Table of contents}
{
\setcounter{tocdepth}{3}
\tableofcontents
}

\section{FUSION Update Documentation
Suite}\label{fusion-update-documentation-suite}

Welcome to the technical documentation suite for \textbf{FUSION Update},
a modern GDAL-based raster handling and native LAS/LAZ point cloud
engine for forestry lidar data processing, terrain analysis, individual
tree detection, and batch processing.

Originally developed by Robert J. McGaughey at the USDA Forest Service /
Pacific Northwest Research Station, the FUSION software suite provides
comprehensive tools for processing airborne laser scanning (ALS) and
terrestrial point cloud datasets. This modern update replaces legacy
binary formats with standard GDAL rasters (GeoTIFF, VRT) and direct
\texttt{.las}/\texttt{.laz} reader/writer infrastructure.

\begin{center}\rule{0.5\linewidth}{0.5pt}\end{center}

\subsection{Documentation Contents}\label{documentation-contents}

This documentation suite is organized into three primary reference
guides:

\subsubsection{\texorpdfstring{1. \href{API_REFERENCE.md}{C++ Core
Library API
Reference}}{1. C++ Core Library API Reference}}\label{c-core-library-api-reference}

Detailed specification of the \texttt{libfusion\_core} developer API,
including headers, namespaces, classes, methods, signatures, parameters,
return types, and memory management patterns: -
\textbf{\texttt{fusion::lidar::LASReader} \& \texttt{LASWriter}}: Point
cloud reader/writer supporting LAS 1.0--1.4 and LAZ compression. -
\textbf{\texttt{fusion::raster::GDALRaster}}: Multi-band and single-band
raster processing, elevation interpolation, and VRT compilation. -
\textbf{\texttt{fusion::batch::BatchPipeline}}: Parallel multi-threaded
spatial tiling engine with buffer management. -
\textbf{\texttt{fusion::metrics}}: High-performance experimental lidar
metric computation engine. -
\textbf{\texttt{fusion::cli::ArgumentParser}}: Command-line argument
parsing and flag handling. -
\textbf{\texttt{fusion::batch::StatusMessenger}}: Thread-safe status
logging and progress monitor.

\subsubsection{\texorpdfstring{2. \href{METRICS_REFERENCE.md}{Lidar \&
Terrain Metrics
Reference}}{2. Lidar \& Terrain Metrics Reference}}\label{lidar-terrain-metrics-reference}

Scientific manual detailing every metric computed by
\texttt{gridmetrics}, \texttt{cloudmetrics}, and
\texttt{ExperimentalMetrics}: - Statistical elevation metrics (mean,
median, percentiles, skewness, kurtosis, IQR). - Canopy cover, canopy
relief ratio (\(CRR\)), and height threshold proportions. - Return
intensity statistics and pulse/return density ratios. - Novel 2D and 3D
spatial metrics (relative height ratios, radial distance variance,
XY/XYZ correlations, columnar grid volume, and 3D voxel volume
occupancy).

\subsubsection{\texorpdfstring{3.
\href{CLI_TOOLS_REFERENCE.md}{Executable Suite \& CLI Tools
Reference}}{3. Executable Suite \& CLI Tools Reference}}\label{executable-suite-cli-tools-reference}

User guide for all 13 command-line executables: -
\texttt{gridmetrics.exe}, \texttt{cloudmetrics.exe},
\texttt{canopymodel.exe}, \texttt{canopymaxima.exe},
\texttt{treeseg.exe}, \texttt{groundfilter.exe}, \texttt{clipdata.exe},
\texttt{filterdata.exe}, \texttt{thindata.exe},
\texttt{returndensity.exe}, \texttt{topometrics.exe},
\texttt{catalog.exe}, and \texttt{ltktools.exe}.

\begin{center}\rule{0.5\linewidth}{0.5pt}\end{center}

\subsection{Architectural Overview}\label{architectural-overview}

\begin{Shaded}
\begin{Highlighting}[]
\NormalTok{graph TD}
\NormalTok{    A["Input Lidar (.las / .laz)"] {-}{-}\textgreater{} B["libfusion\_core"]}
\NormalTok{    C["Input Rasters / DEMs"] {-}{-}\textgreater{} B}
\NormalTok{    B {-}{-}\textgreater{} D["LASPointCloud Engine"]}
\NormalTok{    B {-}{-}\textgreater{} E["GDALRaster Engine"]}
\NormalTok{    B {-}{-}\textgreater{} F["Metrics Calculator"]}
\NormalTok{    B {-}{-}\textgreater{} G["BatchPipeline Engine"]}
\NormalTok{    D {-}{-}\textgreater{} H["13 Executable CLI Tools"]}
\NormalTok{    E {-}{-}\textgreater{} H}
\NormalTok{    F {-}{-}\textgreater{} H}
\NormalTok{    G {-}{-}\textgreater{} H}
\NormalTok{    H {-}{-}\textgreater{} I["Output GeoTIFF Rasters / CSVs"]}
\end{Highlighting}
\end{Shaded}

\begin{center}\rule{0.5\linewidth}{0.5pt}\end{center}

\subsection{Building Documentation}\label{building-documentation}

To re-render this documentation locally into HTML and PDF formats using
Quarto, run the main build script from the repository root:

\begin{Shaded}
\begin{Highlighting}[]
\OperatorTok{.}\NormalTok{\textbackslash{}build}\OperatorTok{.}\FunctionTok{ps1}
\end{Highlighting}
\end{Shaded}

Or execute Quarto directly:

\begin{Shaded}
\begin{Highlighting}[]
\ExtensionTok{quarto}\NormalTok{ render docs}
\end{Highlighting}
\end{Shaded}

The rendered outputs will be placed in: - \texttt{docs/output/html/}:
Interactive HTML documentation site. -
\texttt{docs/pdf/FUSION\_Documentation.pdf}: Standalone PDF reference
manual.




\end{document}
